\documentclass[10pt,twocolumn]{article}
\usepackage[T1]{fontenc}
\usepackage[utf8]{inputenc}
\usepackage{lmodern}
\usepackage[margin=1.8cm,top=2.4cm,bottom=2cm,columnsep=0.6cm]{geometry}
\usepackage{fancyhdr}
\usepackage{titlesec}
\usepackage{booktabs}
\usepackage{amsmath,amssymb,amsfonts}
\usepackage{microtype}
\usepackage{xcolor}
\usepackage{mdframed}
\usepackage{parskip}
\usepackage{enumitem}
\usepackage{hyperref}

\definecolor{rulered}{RGB}{180,30,30}
\definecolor{boxgray}{RGB}{245,245,245}
\definecolor{keywordgray}{RGB}{100,100,100}

\pagestyle{fancy}
\fancyhf{}
\fancyhead[L]{\textsc{\small Claude Mag}}
\fancyhead[C]{\small\textit{Machine Learning Research Digest}}
\fancyhead[R]{\small 2026-09-28}
\fancyfoot[C]{\small\thepage}
\renewcommand{\headrulewidth}{0.8pt}

\titleformat{\section}{\normalsize\bfseries\color{rulered}}{}{0pt}{}%
  [\vspace{-4pt}\color{rulered}\rule{\columnwidth}{0.4pt}]
\titlespacing{\section}{0pt}{8pt}{4pt}

\hypersetup{colorlinks=true,urlcolor=rulered,linkcolor=black}

\begin{document}
\setlength{\parindent}{0pt}
\setlength{\parskip}{4pt}

{\small\color{keywordgray}\textbf{Keywords:}
  test-time compute \textbullet{}
  search policy \textbullet{}
  reinforcement learning \textbullet{}
  LLM reasoning}\par\vspace{4pt}

{\LARGE\bfseries\raggedright
  Beyond Repeated Sampling}\par\vspace{4pt}

{\large\itshape\raggedright
  Learning Search Policies for LLM Reasoning via
  Reinforcement-Trained Concept Generation}\par\vspace{6pt}

{\small\textbf{By} Ismail Labiad, Matthieu Kowalski, Marc Schoenauer,
  R\'emi Munos \& Julia Kempe (Meta FAIR)}\par\vspace{2pt}

{\small\color{keywordgray}arXiv:2609.26704 \textbullet{} September 2026}
\par\vspace{2pt}\noindent{\color{rulered}\rule{\columnwidth}{0.6pt}}\vspace{6pt}

A small concept generator, trained with reinforcement learning to maximize
the pass@k of a frozen large answer generator, achieves substantially higher
diversity and downstream accuracy than naive repeated sampling at the same
compute budget---and transfers to answer generators it was never trained
against, including those from a different model family.

\begin{mdframed}[backgroundcolor=boxgray,linecolor=rulered,linewidth=1pt,
    innerleftmargin=6pt,innerrightmargin=6pt,
    innertopmargin=6pt,innerbottommargin=6pt]
\textbf{\small AT A GLANCE}\par\vspace{4pt}
\begin{description}[leftmargin=0pt,labelsep=4pt,itemsep=2pt,parsep=0pt,
    font=\small\bfseries]
  \item[Problem:] Naive repeated sampling for test-time compute explores only
    through local decoding noise, producing near-duplicate attempts rather
    than semantically diverse solutions.
  \item[Why it matters:] Hard reasoning problems require qualitatively
    different strategies; pass@k saturates quickly when all samples share the
    same implicit approach.
  \item[Method:] Train a small concept generator with RL to emit problem-specific
    hints and strategies that maximize downstream success of a frozen large
    answer generator.
  \item[Result:] Trained concept generator substantially outperforms repeated
    sampling and untuned large models; transfers to unseen answer generators.
  \item[Evidence:] Hard mathematical reasoning benchmarks at equal answer-
    generation compute budget.
\end{description}
\end{mdframed}

\section{The Big Picture}

The dominant paradigm for allocating test-time compute to large language models
is naive repeated sampling: draw $k$ independent answers from the model at
high temperature, then select the best by a verifier or majority vote. This
strategy is seductively simple and, at moderate $k$, effective. But it has a
fundamental limitation: each sample explores only through local decoding
noise. Temperature and top-p diversify token choices within a trajectory,
but they cannot steer the model toward qualitatively different approaches to
the same problem.

On hard mathematical reasoning tasks---where correct solutions often require
an insight, decomposition, or analogy that is far from the model's prior
mode---repeated sampling saturates quickly. The $k$-th sample is nearly as
redundant as the second: the model keeps attempting the same dominant
strategy with minor surface variations.

This paper asks whether test-time exploration can instead be steered at a
\emph{semantic} level. The answer is yes, and the key is to separate the role
of \emph{exploration} (finding diverse approaches) from \emph{execution}
(implementing a given approach): a small, trainable concept generator proposes
strategies; a large, frozen answer generator executes them.

\section{Why Existing Approaches Fall Short}

\textbf{Repeated sampling} (best-of-$N$, majority vote) explores only through
decoding stochasticity. Diversity is local; semantically distinct strategies
require explicit steering.

\textbf{Self-consistency} aggregates repeated samples but does not improve
the diversity of those samples---it is a better aggregator, not a better
explorer.

\textbf{Tree search (MCTS)} explores explicitly but requires a well-trained
per-step verifier, which is expensive to obtain and domain-specific.

\textbf{Prompting for diversity} (e.g., asking the model to use different
methods) provides one alternative per call; emitting a library of diverse
concepts in a single trajectory is more efficient.

\textbf{Scaling the answer generator} at inference time is expensive;
the goal is to make a fixed frozen model explore more effectively with a
lightweight companion.

\section{Core Mechanism}

\textbf{Concept-Conditioned Generation.} Given a problem $x$, the method
first samples a \emph{concept} $c$---a natural-language hint, strategy
keyword, or approach description---then generates an answer from the large
model conditioned on $(x, c)$. The concept shifts the answer generator
toward a different region of solution space.

\textbf{Multi-Concept Trajectories.} Rather than one concept per call, the
concept generator emits a diverse list of concepts in a single forward pass,
covering different strategies before any answer is generated. This maximizes
spread at minimal additional cost.

\textbf{Reinforcement-Learned Concept Generator.} A small model $\pi_\theta$
is trained with RL to emit concepts that maximize the downstream binary
reward (answer correctness). The reward flows back to the concept generator
but not to the frozen answer generator $f$.

\section{Technical Formulation}

Let $x$ denote a problem, $c$ a concept sampled from $\pi_\theta(\cdot \mid x)$,
and $f$ the frozen large answer generator. The concept-conditioned pass@k:
\[
\text{pass@k}(x) = 1 - \frac{\binom{n-s}{k}}{\binom{n}{k}}
\]
where $n$ total samples are generated and $s$ are correct. The RL objective:
\[
\max_\theta \;\mathbb{E}_{x \sim \mathcal{D}}\!\left[
  \text{pass@k}_{c \sim \pi_\theta(\cdot|x)}(f(x, \cdot))
\right]
\]

The policy gradient update uses binary correctness as the scalar reward:
\[
\nabla_\theta \mathcal{L}(\theta) = -\mathbb{E}_{c \sim \pi_\theta}\!\left[
  R(x, c) \, \nabla_\theta \log \pi_\theta(c \mid x)
\right]
\]
where $R(x, c) = \mathbf{1}[f(x,c) \text{ is correct}]$ (or a normalized
group-relative variant for variance reduction).

\section{Learning or Inference Procedure}

\textbf{Training:}
\begin{enumerate}[itemsep=2pt,parsep=0pt]
  \item Initialize $\pi_\theta$ from a small pretrained LM.
  \item For each training problem $x$: sample $K$ concepts from $\pi_\theta$,
    query $f$ for each, record correctness.
  \item Update $\pi_\theta$ with policy gradient on the correctness reward.
\end{enumerate}

\textbf{Inference:}
\begin{enumerate}[itemsep=2pt,parsep=0pt]
  \item Sample a batch of concepts $c_1,\ldots,c_K \sim \pi_\theta(\cdot \mid x)$.
  \item Generate answers $a_i = f(x, c_i)$ for each concept.
  \item Return the most frequent or highest-confidence answer.
\end{enumerate}

\section{What the Guarantee Says}

The trained concept generator is a reusable search policy: it does not learn
to solve problems directly but learns which types of strategies lead the
answer generator to success on the training distribution. Transfer to new
answer generators follows because the learned strategies are problem-structural
rather than model-specific.

\section{Experimental Findings}

\begin{itemize}[itemsep=2pt,parsep=0pt]
  \item Hard mathematical reasoning benchmarks where repeated sampling
    pass@k plateaus at high $k$
  \item Concept-conditioned generation with trained $\pi_\theta$ achieves
    substantially higher pass@k than naive repeated sampling at equal
    answer-generation compute budget
  \item Trained small generator outperforms concepts from untuned models
    many times larger, confirming that RL training---not model scale---is
    the key factor
  \item Zero-shot transfer to an answer generator from a different model
    family succeeds: concepts learned for one model improve another
  \item Multi-concept trajectories (many concepts per call) further improve
    pass@k by reducing concept redundancy
\end{itemize}

\section{Ablations and Interpretation}

\textbf{Single vs.\ multi-concept:} Emitting a diverse list of concepts in
one trajectory reduces redundancy and improves pass@k compared to sampling
one concept per call.

\textbf{Generator size:} Smaller trained generators outperform larger
zero-shot generators, demonstrating the value of training over scale.

\textbf{Reward signal:} Binary correctness alone is sufficient as the RL
reward; no dense process reward model or verifier is needed during training.

\textbf{Transfer:} Concept quality transfers across answer generators from
different model families, suggesting that the search policy captures
problem-structural strategies rather than artifacts of a specific model's
output distribution.

\section*{Reference}

Ismail Labiad, Matthieu Kowalski, Marc Schoenauer, R\'emi Munos, Julia Kempe.
\textbf{Beyond Repeated Sampling: Learning Search Policies for LLM Reasoning}.
arXiv:2609.26704, September 2026.\\
\url{https://arxiv.org/abs/2609.26704}

\end{document}
